% =====================================================================
\chapter{Regras práticas para analisar código}\label{cap:12}
\naaula{46}

Contar cada passo, como no capítulo 2, é ótimo para aprender — mas trabalhoso. Na prática,
reconhecemos alguns \textbf{padrões} e chegamos direto à notação $O$ (ou $\Theta$).

\section{Os padrões}

\begin{center}
\begin{tabularx}{\textwidth}{P{4.3cm} P{5.2cm} L}
\cab \tc{Padrão} & \tc{Exemplo} & \tc{Custo} \\
comandos simples, sem laço & \codigo{x = v[0] + 1} & $\Theta(1)$ \\
\rowcolor{edfundo} um laço de $n$ voltas com corpo $\Theta(1)$ & \codigo{for i in range(n): soma += v[i]} & $\Theta(n)$ \\
laço dentro de laço, $n$ voltas cada & \codigo{for i ... for j ...} & $\Theta(n^2)$ \\
\rowcolor{edfundo} laço \enquote{triangular}: o de dentro vai até \codigo{i} & \codigo{for i in range(n): for j in range(i)} & $\Theta(n^2)$ (Gauss) \\
laço que divide (ou multiplica) por 2 & \codigo{while i > 1: i //= 2} & $\Theta(\log n)$ \\
\rowcolor{edfundo} blocos em sequência & bloco $A$ e depois bloco $B$ & soma: fica o maior \\
\codigo{if}/\codigo{else} & ramo $A$ ou ramo $B$ & o pior dos ramos (para $O$) \\
\rowcolor{edfundo} chamada de função dentro de laço & $n$ chamadas de algo $\Theta(m)$ & $\Theta(n \cdot m)$ \\
\end{tabularx}
\end{center}

\begin{intuicao}
Duas perguntas resolvem quase tudo: \textbf{quantas vezes cada laço roda?} e \textbf{o que está
dentro de quê?} Laços um \emph{depois} do outro somam; laços um \emph{dentro} do outro multiplicam.
\end{intuicao}

\section{Exemplos resolvidos}

\needspace{10\baselineskip}\textbf{Exemplo 1 — laços em sequência e aninhados.}

\begin{lstlisting}
for i in range(n):          # bloco A: n * n voltas
    for j in range(n):
        print(i, j)
for k in range(n):          # bloco B: n voltas
    print(k)
\end{lstlisting}

O bloco A custa $n^2$; o B, $n$. Em sequência, somam: $n^2 + n$. Fica o termo dominante:
$\Theta(n^2)$.

\needspace{10\baselineskip}\textbf{Exemplo 2 — dividindo ao meio.}

\begin{lstlisting}
i = n
while i > 1:
    i = i // 2
\end{lstlisting}

A cada volta, \codigo{i} cai pela metade: o laço roda cerca de $\log_2 n$ vezes. Custo:
$\Theta(\log n)$.

\needspace{10\baselineskip}\textbf{Exemplo 3 — o laço triangular.}

\begin{lstlisting}
for i in range(n):
    for j in range(i):      # 0, 1, 2, ..., n-1 voltas
        print(i, j)
\end{lstlisting}

O laço de dentro dá $0 + 1 + 2 + \dots + (n - 1) = \dfrac{n(n-1)}{2}$ voltas no total (Gauss).
Custo: $\Theta(n^2)$ — a metade de um quadrado continua quadrática.

\needspace{10\baselineskip}\textbf{Exemplo 4 — um laço logarítmico dentro de um linear.}

\begin{lstlisting}
for i in range(n):          # n voltas
    j = 1
    while j < n:            # cerca de log2(n) voltas
        j = j * 2
\end{lstlisting}

Para cada uma das $n$ voltas de fora, o de dentro roda cerca de $\log_2 n$ vezes. Aninhados,
multiplicam: $\Theta(n \log n)$.

\needspace{10\baselineskip}\textbf{Exemplo 5 — a função escondida.}

\begin{lstlisting}
def comuns(a, b):           # a tem n elementos, b tem m
    resultado = []
    for x in a:             # n voltas
        if x in b:          # "in" percorre b: O(m)
            resultado.append(x)
    return resultado
\end{lstlisting}

O \codigo{x in b} parece um teste simples, mas é uma busca linear em \codigo{b}
(capítulo 2). São $n$ buscas de custo $m$: $O(n \cdot m)$. É a interseção ingênua da Aula 3.

\begin{atencao}
Em Python, desconfie de toda linha que opera sobre uma lista inteira: \codigo{in},
\codigo{insert(0, x)}, \codigo{remove(x)}, \codigo{index(x)}, \codigo{min}, \codigo{max},
\codigo{sum}, \codigo{sorted}. Cada uma esconde um laço.
\end{atencao}

\begin{exrapido}[14]
Dê o custo de cada trecho, em notação $O$.

\begin{lstlisting}
# (a)
for i in range(n):
    for j in range(10):
        print(i, j)
# (b)
i = 1
while i < n:
    i = i * 3
# (c)
for i in range(n):
    lista.insert(0, i)      # lista contígua do Python
\end{lstlisting}
\end{exrapido}

\begin{respostas}
  \resp{12.1} (a) O laço de dentro tem sempre 10 voltas (uma constante): $10n$, ou seja, $O(n)$.
  (b) \codigo{i} é multiplicado por 3 a cada volta: $\log_3 n$ voltas, ou seja, $O(\log n)$ (a
  base não importa). (c) Cada \codigo{insert(0, i)} desloca todos os elementos já inseridos:
  $0 + 1 + \dots + (n - 1)$ deslocamentos, ou seja, $O(n^2)$.
\end{respostas}
